% 文件开头添加条件判断
\newif\ifstandalone
\standalonetrue  % 默认独立编译模式

% 检测是否在主文档中
\ifdefined\maindoc
  \standalonefalse
\fi

\ifstandalone
  % 第四章 计算力矩控制
  \documentclass[UTF8]{ctexart}

  % 页面设置
  \usepackage[a4paper, margin=2.5cm]{geometry}
  \usepackage{setspace}
  \onehalfspacing

  % 数学包
  \usepackage{amsmath}
  \usepackage{amssymb}
  \usepackage{amsthm}
  \usepackage{mathrsfs}
  \usepackage{bm}

  % 图形包
  \usepackage{tikz}
  \usetikzlibrary{shapes,arrows,positioning,calc}
  \usepackage{graphicx}

  % 表格包
  \usepackage{booktabs}
  \usepackage{array}
  \usepackage{multirow}
\fi

\ifstandalone
% 列表包
\usepackage{enumitem}
\usepackage{listings}

% 其他包
\usepackage{float}
\usepackage{caption}
\usepackage{subcaption}
\usepackage{hyperref}
\usepackage{xcolor}

% 定义定理环境
\theoremstyle{definition}
\newtheorem{theorem}{定理}[section]
\newtheorem{lemma}{引理}[section]
\newtheorem{example}{示例}[section]
\newtheorem{remark}{注记}[section]

% 定义颜色
\definecolor{examplebg}{RGB}{245,245,250}
\definecolor{exampleframe}{RGB}{100,100,120}

% 示例环境
\usepackage{tcolorbox}
\tcbuselibrary{skins,breakable}
\newtcolorbox{examplebox}[1][]{
    colback=examplebg,
    colframe=exampleframe,
    fonttitle=\bfseries,
    title={#1},
    breakable,
    enhanced
}

% 标题信息
\title{\textbf{第四章 计算力矩控制}\\
\large Computed-Torque Control}
\author{机器人操作臂控制——理论与实践}
\date{}

% 自定义命令
\newcommand{\vect}[1]{\boldsymbol{#1}}
\newcommand{\mat}[1]{\boldsymbol{#1}}
\newcommand{\trans}{\mathrm{T}}
\newcommand{\dd}{\mathrm{d}}
\newcommand{\ATAN}{\mathrm{ATAN2}}
\fi

\ifstandalone
\begin{document}
\maketitle
\fi

\section{引言}

控制机器人的一个基本问题是使操作臂跟随预先规划的期望轨迹。在机器人执行任何有用的工作之前，我们必须在正确的时刻将其定位在正确的位置。本章讨论计算力矩控制，它产生一系列易于理解的控制方案，在实践中通常效果很好。这些方案涉及将控制设计问题分解为内环设计和外环设计。

在4.4节中，我们将提供基于PID控制的独立关节设计的经典操作臂控制方案的联系。在4.6节中，我们将展示如何结合计算力矩控制使用一些现代设计技术。因此，本章可以被视为连接几年前机器人控制中使用的经典设计技术与本书其余部分中需要的现代设计技术之间的桥梁，以在不确定环境中获得高性能。

这里我们假设机器人在自由空间中移动，不与环境接触。接触会导致力的产生。力控制问题在第7章中讨论。我们还将假设在本章中机器人是一个已知的刚体系统，因此基于一个相当已知的模型设计控制器。在存在不确定性或未知参数（例如摩擦、负载质量）的情况下进行控制需要精细的方法。这个问题在第4章使用鲁棒控制处理，在第5章使用自适应控制处理。

实际的机器人操作臂可能在其连杆中具有柔性，或在其齿轮中具有顺应性（关节柔性）。在第6章中，我们涵盖了带关节柔性的控制的一些方面。

在我们可以控制机器人手臂之前，有必要知道执行任务的期望路径。与路径规划问题相关的有许多问题，例如避免障碍物以及确保规划的路径不需要超过执行器的电压和力矩限制。为了将控制问题简化为基本组成部分，在本章中，我们假设最终的控制目标是使机器人沿着规定的期望轨迹移动。我们不关心实际的轨迹规划问题；但是，我们确实展示了如何重建一个连续的期望路径，该路径从给定期望点表（末端执行器应通过这些点）中得出。这个连续路径生成问题在4.2节中介绍。

在大多数情况下，机器人控制器在微处理器上实现，特别是考虑到现代控制方案的复杂性质。因此，在4.5节中，我们说明了机器人控制器数字实现的一些概念。

贯穿始终，我们演示如何在计算机上模拟机器人控制器。在实际在真实机器人操作臂上实现之前，应该这样做以验证任何提议的控制方案的有效性。

\section{路径生成}

在本书中，我们假设给定了一条规定路径$q_d(t)$，机器人手臂应该跟随它。我们设计控制方案使操作臂跟随这个期望路径或轨迹。轨迹规划涉及找到规定路径，通常被视为一个单独的设计问题，涉及碰撞避免、执行器饱和等考虑。参见[Lee et al. 1983]。

我们不涵盖轨迹规划。但是，我们确实涵盖轨迹生成的两个方面。首先，我们展示了如何将给定的规定路径从笛卡尔空间转换到关节空间。然后，给定一个期望点表（末端执行器应通过这些点），我们展示了如何重建一个连续的期望轨迹。

\subsection{将笛卡尔轨迹转换到关节空间}

在机器人应用中，期望任务通常在工作空间或笛卡尔空间中指定，因为在这里操作臂的运动很容易相对于外部环境和工作件进行描述。然而，轨迹跟随控制在关节空间中很容易执行，因为这里是臂动力学更容易表述的地方。

因此，重要的是能够找到给定期望笛卡尔轨迹的期望关节空间轨迹$q_d(t)$)。这通过使用逆运动学来完成，如下面示例所示。该示例说明了笛卡尔到关节空间轨迹的映射可能不是唯一的——也就是说，几个关节空间轨迹可能对末端执行器产生相同的笛卡尔轨迹。

\begin{examplebox}[示例4.2-1：将规定笛卡尔轨迹映射到关节空间]
在示例A.3-5中推导出了图4.2.1所示的两连杆平面机器人手臂的逆运动学。让我们使用它们将路径从笛卡尔空间转换到关节空间。

假设我们希望两连杆臂在$(x, y)$平面中跟随给定的工作空间或笛卡尔轨迹
\begin{equation}
    p(t) = (x(t), y(t))
\end{equation}
它是时间$t$的函数。由于手臂由控制其角度$\theta_1, \theta_2$的执行器移动，因此方便将指定的笛卡尔轨迹$(x(t), y(t))$转换为关节空间轨迹$(\theta_1(t), \theta_2(t))$以进行控制。

这可以通过使用逆运动学变换来实现
\begin{align}
    r^2 &= x^2 + y^2 \\
    C &= \cos\theta_2 = \frac{r^2 - a_1^2 - a_2^2}{2a_1a_2} \\
    D &= \pm\sqrt{1 - C^2} = \sin\theta_2 \\
    \theta_2 &= \ATAN(D, C) \\
    \theta_1 &= \ATAN(y, x) - \ATAN(a_2\sin\theta_2, a_1 + a_2\cos\theta_2)
\end{align}

假设臂的末端应重复追踪图4.2.2所示的圆形工作空间路径$p(t)$，其描述为
\begin{align}
    x(t) &= 2 + \frac{1}{2}\cos t \\
    y(t) &= 1 + \frac{1}{2}\sin t
\end{align}

通过在每个时间$t$在逆运动学方程中使用这些表达式，我们得到图4.2.3所示的所需关节空间轨迹$q(t) = (\theta_1(t), \theta_2(t))$，这些轨迹产生了末端执行器的圆形笛卡尔运动（使用$a_1 = 2, a_2 = 2$）。

我们为"肘部向下"配置计算了关节变量。在(4)中选择相反的符号会得到"肘部向上"的关节空间轨迹，产生相同的笛卡尔轨迹。
\end{examplebox}

\subsection{多项式路径插值}

假设已经为操作臂运动确定了期望轨迹，无论是在笛卡尔空间还是使用逆运动学在关节空间。为方便起见，我们使用关节空间变量$q(t)$表示。将整个轨迹存储在计算机内存中是不可能的，很少有实用有用的轨迹具有简单的闭合形式表达式。因此，通常的做法是将一系列点$q_i(t_k)$存储在计算机内存中，对于每个关节变量$i$，这些点代表变量在离散时间$t_k$的期望值。因此$q(t_k)$是$\mathbb{R}^n$中的一个点，关节变量应在时间$t_k$通过该点。我们称这些为经过点。

大多数机器人控制方案需要连续的期望轨迹。为了将经过点表$q_i(t_k)$转换为连续的期望轨迹$q_d(t)$，我们可以使用许多选项。让我们在这里讨论多项式插值。

假设经过点在时间上是均匀间隔的，并将采样周期定义为
\begin{equation}
    T = t_{k+1} - t_k
\end{equation}

对于平滑运动，在每个时间间隔$[t_k, t_{k+1}]$上，我们要求期望位置$q_d(t)$和速度$\dot{q}_d(t)$与制表的经过点匹配。这产生了边界条件
\begin{equation}
    \begin{cases}
        q_d(t_k) = q(t_k), & q_d(t_{k+1}) = q(t_{k+1}) \\
        \dot{q}_d(t_k) = \dot{q}(t_k), & \dot{q}_d(t_{k+1}) = \dot{q}(t_{k+1})
    \end{cases}
\end{equation}

为了匹配这些边界条件，有必要在每个时间间隔$[t_k, t_{k+1}]$上使用三次插值多项式
\begin{equation}
    q_d(t) = a_0 + a_1t + a_2t^2 + a_3t^3
\end{equation}
它有四个自由变量。那么
\begin{align}
    \dot{q}_d(t) &= a_1 + 2a_2t + 3a_3t^2 \\
    \ddot{q}_d(t) &= 2a_2 + 6a_3t
\end{align}

因此，加速度在每个采样周期上是线性的。

很容易求解在每个时间间隔$[t_k, t_{k+1}]$上保证边界条件匹配的系数。实际上，我们看到
\begin{equation}
    \begin{bmatrix} 1 & t_k & t_k^2 & t_k^3 \\ 0 & 1 & 2t_k & 3t_k^2 \\ 1 & t_{k+1} & t_{k+1}^2 & t_{k+1}^3 \\ 0 & 1 & 2t_{k+1} & 3t_{k+1}^2 \end{bmatrix} \begin{bmatrix} a_0 \\ a_1 \\ a_2 \\ a_3 \end{bmatrix} = \begin{bmatrix} q(t_k) \\ \dot{q}(t_k) \\ q(t_{k+1}) \\ \dot{q}(t_{k+1}) \end{bmatrix}
\end{equation}

这被求解以获得每个时间间隔$[t_k, t_{k+1}]$上所需的插值系数。

请注意，这种技术需要以表格形式存储每个采样点的期望位置和速度。一种变体使用更高阶的多项式来确保在每个采样时间$t_k$处的位置、速度和加速度连续。

虽然我们已经使用了关节变量符号$q(t)$，但应该强调的是，轨迹插值也可以在笛卡尔空间中进行。

\subsection{带抛物线混合的线性函数}

使用三次插值多项式，加速度在每个采样周期上是线性的。然而，在许多实际应用中，有很好的理由坚持在每个采样周期内保持恒定加速度。例如，任何真实的机器人都有一个执行器可以提供的力矩的上限。对于线性系统（想想牛顿定律），这转化为恒定加速度的上限。因此，恒定加速度不太可能导致执行器饱和。除此之外，大多数工业机器人控制器被编程为在每个采样周期上使用恒定加速度。

恒定加速度曲线如图4.2.4(a)所示。相关的速度和位置曲线如图4.2.4(b)和4.2.4(c)所示。位置轨迹由三部分组成：二次或抛物线初始部分、线性中间部分和抛物线最终部分。因此，让我们讨论使用带抛物线混合的线性函数(LFPB)的经过点插值。

位置轨迹从抛物线切换到线性的时间被称为混合时间$t_b$。应为每个关节变量$i$指定位置$q_{di}(t)$。图4.2.4(c)中的轨迹可以写成关节$i$的形式
\begin{equation}
    q_{di}(t) = \begin{cases}
        a_{i0} + a_{i1}t + a_{i2}t^2 & 0 \leq t < t_b \\
        b_{i0} + b_{i1}t & t_b \leq t < t_f - t_b \\
        c_{i0} + c_{i1}t + c_{i2}t^2 & t_f - t_b \leq t \leq t_f
    \end{cases}
\end{equation}

系数$v_i$可以被解释为关节变量$i$允许的最大速度。设计参数是$v_i$和$t_b$。

很容易求解在每个时间间隔$[t_k, t_{k+1}]$上确保满足边界条件(4.2.2)的系数。结果是
\begin{equation}
\left\{
\begin{aligned}
        a_{i0} &= q_{ik}, & a_{i1} &= v_i/t_b, & a_{i2} &= -v_i/(2t_b) \\
        b_{i0} &= q_{ik} + v_it_b/2, & b_{i1} &= v_i \\
        c_{i0} &= q_{i,k+1} - v_it_f^2/(2t_b), & c_{i1} &= v_it_f/t_b, & c_{i2} &= -v_i/(2t_b)
\end{aligned}
\right.
\end{equation}

\subsection{最小时间轨迹}

有一类重要的特殊LFPB轨迹。假设加速度受到最大值$a_M$的限制，并且希望机器人手臂以最小时间从一个位置到达另一个位置。为简单起见，假设初始和最终速度等于零。一般情况在[Lewis 1986a]中涵盖（参见习题）。

最小时间轨迹如图4.2.5所示。为了在最小时间$t_f$内将关节变量$i$从静止位置$q_0 = q_i(t_0)$驱动到期望的最终静止位置$q_f = q_i(t_f)$，最大加速度$a_M$应施加到切换时间$t_s$，之后时间应施加最大减速度$-a_M$直到$t_f$。注意$t_s$和$t_f$都取决于$q_0$和$q_f$。我们可以写成
\begin{equation}
    q_i(t_s) = q_0 + \frac{1}{2}a_M(t_s - t_0)^2
\end{equation}

那么速度方程产生
\begin{equation}
    v_i(t_s) = a_M(t_s - t_0) = a_M(t_f - t_s)
\end{equation}
或
\begin{equation}
    t_s = \frac{t_f + t_0}{2}
\end{equation}

也就是说，从最大加速度到最大减速度的切换发生在半时间点。现在对位置方程进行简单的操作产生
\begin{equation}
    q_f = q_i(t_s) + v_i(t_s)(t_f - t_s) - \frac{1}{2}a_M(t_f - t_s)^2
\end{equation}

因此(4.2.10)产生
\begin{equation}
    t_f = t_0 + 2\sqrt{\frac{q_f - q_0}{a_M}}
\end{equation}

不幸的是，使用恒定最大加速度计算的最小时间轨迹与机器人技术没有直接相关性。这是由于实际的操作臂具有$\tau_m$的力矩限制这一事实。由于机器人方程（见表3.3.1）是非线性的，这与加速度的上限恒定没有相关性。例如，机器人手臂在其完全伸展和完全缩回位置具有不同的最大加速度。有关更多讨论，请参见[Kahn and Roth 1971]、[Chen 1989]、[Geering 1986]、[Gourdeau and Schwartz 1989]、[Jayasuriya and Suh 1985]、[Kim and Shin 1985]、[Shin and McKay 1985]。

\section{机器人系统的计算机仿真}

在实际在臂上实现之前，在数字计算机上仿真提议的操作臂控制方案非常重要。我们在这里展示如何对机器人系统进行这样的计算机仿真。由于大多数机器人控制器实际上以数字方式实现（第4.5节），我们还展示了如何仿真数字机器人手臂控制器。

\subsection{机器人动力学仿真}

有许多用于仿真非线性动力学系统的软件包，包括SIMNON [Åström and Wittenmark 1984]、MATLAB等。为方便起见，我们在附录B中包含了一些对连续和数字控制相当有用的仿真程序。

所有仿真程序都要求用户编写类似的子程序。使用积分程序（如Runge-Kutta）的时间响应仿真器都需要计算给定当前状态的状态导数。在3.4节中，我们展示了如何将机器人手臂方程
\begin{equation}
    M(q)\ddot{q} + N(q, \dot{q}) = \tau
\end{equation}
以非线性状态空间形式表示
\begin{equation}
    \dot{x} = f(x, u)
\end{equation}
其中$x(t)$是状态，$u(t)$是输入。

定义状态为$x = [q^T \, \dot{q}^T]^T$，我们可以写出隐式形式
\begin{equation}
    \begin{bmatrix} I & 0 \\ 0 & M(q) \end{bmatrix} \dot{x} = \begin{bmatrix} 0 & I \\ 0 & -N(q, \dot{q}) \end{bmatrix} x + \begin{bmatrix} 0 \\ I \end{bmatrix} \tau + \begin{bmatrix} 0 \\ I \end{bmatrix} \tau_d
\end{equation}
其中$\tau$是臂控制力矩，由控制器提供，$\tau_d$是干扰力矩。我们说这是一种隐式形式，因为左侧的系数矩阵意味着$\ddot{q}$不是以右侧的显式形式给出的。

给定$x(t)$，有必要为积分程序提供一个子程序来计算$\dot{x}(t)$。一种求解$\ddot{q}$的方法是求逆$M(q)$。然而，由于潜在的数值问题，这不推荐。让我们将(4.3.3)表示为
\begin{equation}
    \dot{x} = g(x, u)
\end{equation}
请注意，在这种情况下，$u(t)$是由控制力$\tau(t)$和干扰$\tau_d(t)$组成的向量。

附录B中给出了一个简单的时间响应程序TRESP。给定一个子程序F(time, x, xp)，它使用(4.3.4)根据$x(t)$和$u(t)$计算$\dot{x}$；它使用Runge-Kutta积分器来计算状态轨迹$x(t)$。为了在子程序F(t, x, xp)中求解$\dot{x}$，我们建议计算$M(q)$和$N(q, \dot{q})$，然后使用最小二乘技术求解
\begin{equation}
    \ddot{q} = M^{-1}(q)[-N(q, \dot{q}) + \tau]
\end{equation}
[i.e., (4.3.3)的底部部分]，这比$M(q)$的逆数值上更稳定。最小二乘方程求解器在商业上很容易获得，例如在[IMSL]、[LINPACK]等处。对于更简单的臂，可以解析地求逆$M(q)$。

贯穿全书，我们使用各种控制方案说明臂动力学的仿真。

\subsection{数字机器人控制器的仿真}

虽然大多数机器人控制器是在连续时间中设计的，但它们在实际的机器人上以数字方式实现。也就是说，控制信号仅在离散的时间瞬间使用微处理器更新。我们在4.5节中讨论数字机器人手臂控制器的实现。因此，在实际实现之前验证提议的控制器是否按预期运行是非常可取的，使用其数字化或离散形式对其进行仿真。

数字控制方案如图4.3.1所示。要控制的工厂或系统是连续时间系统，$K(z)$是动态数字控制器，其中$z$是Z变换变量（即，$z^{-1}$代表单位时间延迟）。数字控制器$K(z)$使用DSP中的软件代码实现。参考输入$r(t)$是$y(t)$应该跟随的期望轨迹，$e_k$是（离散）跟踪误差。

采样器与采样周期$T$是一个模数(A/D)转换器，它获取软件控制器$K(z)$所需的输出$y(t)$的样本$y_k = y(kT)$。$y(t)$的定义可以根据控制方案而变化。例如，在机器人控制中，$y(t)$可能表示由$q(t)$和$\dot{q}(t)$组成的$2n$向量。

保持设备是图中的一个数模(D/A)转换器，它将软件控制器$K(z)$计算的离散控制样本$u_k$转换为工厂所需的连续时间控制$u(t)$。它是一种数据重建设备。零阶保持(ZOH)的输入$u_k$和输出$u(t)$如图4.3.2所示。注意$u(kT) = u_k$，$T$为采样周期，因此$u(t)$从右侧是连续的。也就是说，$u(t)$在时刻$kT$更新。ZOH通常用于控制目的，而不是其他高阶设备如一阶保持，因为大多数商用DSP都有内置的ZOH。

一旦设计了控制器，重要的是在使用数字计算机实现之前对其进行仿真，以确定闭环响应是否合适。这在机器人技术中尤其正确，因为数字控制器通常是通过设计连续时间控制器找到的，然后使用Euler方法等近似技术进行数字化。也就是说，对于非线性系统，控制器离散化方案通常不是精确的。这会导致性能下降。为了验证控制器性能是否合适，仿真应在所有时间提供响应，包括采样之间的时间。

为了仿真数字控制器，我们可以使用图4.3.3所示的方案。在那里，连续工厂动力学包含在子程序F(t, x, xp)中；它们使用Runge-Kutta积分器进行积分。该图假设ZOH；因此控制输入$u(t)$在每个时间$kT$更新为$u_k$，然后保持恒定直到时间$(k+1)T$。注意涉及两个时间间隔；采样周期$T$和Runge-Kutta积分周期$T_R \ll T$。$T_R$应被选为$T$的整数除数。

这种仿真技术将工厂状态$x(t)$提供为时间的连续函数，即使在采样瞬间之间的值也是如此[事实上，它以$T_R$的倍数提供$x(t)$]。这对于在实际在工厂上实现数字控制器之前验证闭环系统的可接受的样本间行为是至关重要的。

附录B中的程序TRESP可用于实现图4.3.3。它以模块化的方式编写，可应用于广泛的情况。我们将在随后的几个示例中说明其用于数字控制的用途。使用SIMNON等仿真软件非常相似。

我们在4.5节中讨论数字机器人手臂控制器的实现。[Lewis 1992]中给出了关于数字控制、仿真和DSP控制器实现的详细讨论。

\section{计算力矩控制}

多年来提出了许多机器人控制方案。事实上，大多数都可以被视为计算力矩控制器类的特殊情况。同时，计算力矩是非线性系统反馈线性化的特殊应用，这在现代系统理论中很流行[Hunt et al. 1983]、[Gilbert and Ha 1984]。事实上，对机器人控制方案进行分类的一种方法是将其分为"类计算力矩"或"非类计算力矩"。类计算力矩控制出现在鲁棒控制、自适应控制、学习控制等中。

在本章的剩余部分，我们探索这类机器人控制器，它包括广泛的设计。计算力矩控制使我们能够方便地推导出非常有效的机器人控制器，同时提供一个框架，将经典独立关节控制和一些现代设计技术结合在一起，并为本书的其余部分奠定基础。在本节末尾的表4.4.1中总结了不同类型的类计算力矩控制器。我们将看到，许多数字机器人控制器也是类计算力矩控制器（第4.5节）。

\subsection{内前馈环推导}

机器人手臂动力学为
\begin{equation}
    M(q)\ddot{q} + N(q, \dot{q}) = \tau
\end{equation}
或
\begin{equation}
    M(q)\ddot{q} + V_m(q, \dot{q})\dot{q} + G(q) + F(\dot{q}) = \tau - \tau_d
\end{equation}
其中关节变量$q(t)$、控制力矩$\tau$和干扰$\tau_d(t)$。如果此方程包括电机执行器动力学（第3.6节），则$\tau(t)$是输入电压。

假设已根据4.2节的讨论为臂运动选择了期望轨迹$q_d(t)$。为了确保关节变量的轨迹跟踪，定义输出或跟踪误差为
\begin{equation}
    e(t) = q_d(t) - q(t)
\end{equation}

为了演示输入$\tau(t)$对跟踪误差的影响，微分两次以获得
\begin{equation}
    \ddot{e} = \ddot{q}_d - \ddot{q}
\end{equation}

现在求解(4.4.2)中的$\ddot{q}$并将其代入最后一个方程得到
\begin{equation}
    \ddot{e} = \ddot{q}_d - M^{-1}(q)[-N(q, \dot{q}) + \tau - \tau_d]
\end{equation}

定义控制输入函数
\begin{equation}
    u = M^{-1}(q)[-N(q, \dot{q}) + \tau] - \ddot{q}_d
\end{equation}
和干扰函数
\begin{equation}
    w = M^{-1}(q)\tau_d
\end{equation}

我们可以按如下方式定义状态$x = [e^T \, \dot{e}^T]^T$并写出跟踪误差动力学
\begin{equation}
    \frac{\dd}{\dd t}\begin{bmatrix} e \\ \dot{e} \end{bmatrix} = \begin{bmatrix} 0 & I \\ 0 & 0 \end{bmatrix}\begin{bmatrix} e \\ \dot{e} \end{bmatrix} + \begin{bmatrix} 0 \\ I \end{bmatrix}u + \begin{bmatrix} 0 \\ I \end{bmatrix}w
\end{equation}

这是Brunovsky规范形中的线性误差系统，由$n$对双积分器$1/s^2$组成，每个关节一个。它由控制输入$u(t)$和干扰$w(t)$驱动。注意，此推导是3.4节中一般反馈线性化程序的特殊情况。

反馈线性化变换(4.4.5)可以反转以得到
\begin{equation}
    \tau = M(q)(\ddot{q}_d - u) + N(q, \dot{q})
\end{equation}

我们称之为计算力矩控制律。这些操作的重要性如下。从(4.4.1)到(4.4.8)没有状态空间变换。因此，如果我们选择一个稳定(4.4.8)的控制$u(t)$，使得$e(t)$趋于零，那么由$\tau(t)$(4.4.9)给出的非线性控制输入将导致机器人臂(4.4.1)中的轨迹跟随。事实上，将(4.4.9)代入(4.4.2)得到
\begin{equation}
    M\ddot{q} + N = M(\ddot{q}_d - u) + N + \tau_d
\end{equation}
或
\begin{equation}
    \ddot{e} = u + M^{-1}\tau_d = u + w
\end{equation}
这正是(4.4.8)。

图4.4.1展示了计算力矩控制方案，显示了内环和外环。重要的是要注意，它由一个内非线性环和一个外控制信号$u(t)$组成。我们将看到选择$u(t)$的几种方法。由于$u(t)$将取决于$q(t)$和$\dot{q}(t)$，外环将是一个反馈环。一般来说，我们可以选择一个动态补偿器$H(s)$，使得
\begin{equation}
    U(s) = H(s)E(s)
\end{equation}

$H(s)$可以被选择用于良好的闭环行为。根据(4.4.10)，闭环误差系统然后具有传递函数
\begin{equation}
    T(s) = s^2I - H(s)
\end{equation}

重要的是要意识到计算力矩取决于机器人动力学的反演，实际上有时被称为逆动力学控制。事实上，(4.4.9)显示了通过用$\ddot{q}_d - u$替换$\ddot{q}$来计算$\tau(t)$；也就是说，通过求解机器人逆动力学问题。与系统反演相关的所有警告，包括当系统具有非最小相位零点时产生的问题，都在这里适用。（注意在线性情况下，系统零点是反演的极点。这种非最小相位概念推广到非线性系统。）幸运的是，对于我们来说，刚性臂动力学是最小相位的。

出于实现目的，有几种方法可以计算(4.4.9)。应避免在每个采样时间进行形式矩阵乘法。在某些情况下，表达式可以解析地推导出来。计算力矩$\tau(t)$的一个好方法是使用高效的Newton-Euler逆动力学公式[Craig 1985]，用$\ddot{q}_d - u$代替$\ddot{q}(t)$。

可以使用许多方法选择外环信号$u(t)$，包括鲁棒和自适应控制技术。在本章的剩余部分，我们探索$u(t)$的一些选择以及计算力矩控制的一些变体。

\subsection{PD外环设计}

选择辅助控制信号$u(t)$的一种方法是作为比例加微分(PD)反馈
\begin{equation}
    u = -K_v\dot{e} - K_pe
\end{equation}

然后整体机器人臂输入变为
\begin{equation}
    \tau_c = M(q)(\ddot{q}_d + K_v\dot{e} + K_pe) + N(q, \dot{q})
\end{equation}

此控制器如图4.4.6所示，其中$K_i = 0$。闭环误差动力学为
\begin{equation}
    \ddot{e} + K_v\dot{e} + K_pe = w
\end{equation}
或以状态空间形式
\begin{equation}
    \frac{\dd}{\dd t}\begin{bmatrix} e \\ \dot{e} \end{bmatrix} = \begin{bmatrix} 0 & I \\ -K_p & -K_v \end{bmatrix}\begin{bmatrix} e \\ \dot{e} \end{bmatrix} + \begin{bmatrix} 0 \\ I \end{bmatrix}w
\end{equation}

闭环特征多项式为
\begin{equation}
    \Delta_c(s) = |s^2I + sK_v + K_p|
\end{equation}

\textbf{PD增益选择}。通常取$n \times n$增益矩阵为对角矩阵，使得
\begin{equation}
    K_v = \mathrm{diag}\{k_{vi}\}, \quad K_p = \mathrm{diag}\{k_{pi}\}
\end{equation}

那么
\begin{equation}
    \Delta_c(s) = \prod_{i=1}^{n}\left(s^2 + sk_{vi} + k_{pi}\right)
\end{equation}
只要所有的$k_{vi}$和$k_{pi}$都为正，误差系统就是渐近稳定的。因此，只要干扰$w(t)$有界，误差$e(t)$就是有界的。与此相关，检查(4.4.6)并回忆表3.3.1，$M^{-1}$是有上界的。因此，$w(t)$的有界性等同于$\tau_d(t)$的有界性。

重要的是要注意，虽然选择对角PD增益矩阵导致在外环级别解耦控制，但它不会导致解耦的关节控制策略。这是因为内环中的乘以$M(q)$和加上非线性前馈项$N(q, \dot{q})$使得信号$u(t)$在所有关节之间混杂。因此，计算任何一个给定关节的控制$\tau(t)$通常需要所有关节位置$q(t)$和速度$\dot{q}(t)$的信息。

二阶特征多项式的标准形式为
\begin{equation}
    s^2 + 2\zeta\omega_ns + \omega_n^2
\end{equation}
其中$\zeta$是阻尼比，$\omega_n$是自然频率。因此，可以通过如下选择PD增益来实现误差$e(t)$的每个分量中的期望性能
\begin{equation}
    k_{pi} = \omega_{ni}^2, \quad k_{vi} = 2\zeta_i\omega_{ni}
\end{equation}
其中$\zeta_i, \omega_{ni}$是关节误差$i$的期望阻尼比和自然频率。使末端附近的响应比靠近基座处更快可能是有用的，因为必须移动的质量更重。

机器人表现出超调是不希望的，因为如果例如期望轨迹在工件的表面终止，这可能导致冲击。因此，PD增益通常选择为临界阻尼$\zeta = 1$。在这种情况下
\begin{equation}
    k_{pi} = \omega_{ni}^2, \quad k_{vi} = 2\omega_{ni}
\end{equation}

\textbf{自然频率的选择}。自然频率$\omega_n$控制每个误差分量中的响应速度。对于快速响应，它应该很大，并且根据性能目标进行选择。因此，在选择$\omega_n$时应考虑期望轨迹。

现在讨论选择中的一些其他因素。[Paul 1981]给出了$\omega_n$选择上的一些上限。尽管大多数工业机器人的连杆是重的，但它们可能有一些柔性。假设连杆$i$的第一柔性或共振模式的频率为
\begin{equation}
    \omega_r = \sqrt{k_r/J}
\end{equation}
其中$J$是连杆惯量，$k_r$是连杆刚度。那么，为了避免激发共振模式，我们应该选择$\omega_n < \omega_r/2$。当然，连杆惯量$J$随臂构型变化，因此其最大值可用于计算$\omega_r$。

$\omega_n$上的另一个上限由执行器饱和考虑提供。如果PD增益太大，力矩$\tau(t)$可能达到其上限。

如下从误差有界性考虑提供了PD增益选择的一些更多感觉。(4.4.15)中闭环误差系统的传递函数为
\begin{equation}
    H(s) = (s^2I + sK_v + K_p)^{-1}
\end{equation}
或如果$K_v$和$K_p$是对角的
\begin{equation}
    H(s) = \mathrm{diag}\left\{\frac{1}{s^2 + sk_{vi} + k_{pi}}\right\}
\end{equation}

\begin{equation}
    w = M^{-1}\tau_d + \ddot{q}_d
\end{equation}

我们假设干扰和$M^{-1}$是有界的（表3.3.1），使得
\begin{equation}
    \|w\|_2 \leq \|M^{-1}\|_2\|\tau_d\|_2 + \|\ddot{q}_d\|_2 \leq \bar{m}\|\tau_d\|_2 + \|\ddot{q}_d\|_2 \leq d
\end{equation}
其中对于给定的机器人臂，$m$和$d$是已知的。因此
\begin{equation}
    \|e\|_2 \leq \|H(s)\|_2\|w\|_2 \leq \|H(s)\|_2d
\end{equation}
\begin{equation}
    \|\dot{e}\|_2 \leq \|sH(s)\|_2d
\end{equation}

现在选择$L_2$范数，算子增益$\|H(s)\|_2$是$H(s)$的Bode幅度图的最大值。对于临界阻尼系统
\begin{equation}
    |H_i(j\omega)| = \frac{1}{\sqrt{(\omega^2 - \omega_{ni}^2)^2 + (2\zeta_i\omega_{ni}\omega)^2}}
\end{equation}
因此
\begin{equation}
    \|H(s)\|_2 = \max_i\frac{1}{2\zeta_i\omega_{ni}^2}
\end{equation}

此外（参见习题）
\begin{equation}
    \|sH(s)\|_2 = \max_i\frac{1}{2\zeta_i\omega_{ni}}
\end{equation}
所以
\begin{equation}
    \|e\|_2 \leq \frac{d}{\omega_{ni}^2}, \quad \|\dot{e}\|_2 \leq \frac{d}{\omega_{ni}}
\end{equation}

因此，在临界阻尼的情况下，位置误差随$k_{pi}$减小，速度误差随$k_{vi}$减小。

\begin{examplebox}[示例4.4-1：PD计算力矩控制的仿真]
\textbf{a. 计算力矩控制律}

在示例3.2.2中，我们找到了图4.2.1所示的两连杆平面肘臂的动力学
\begin{equation}
    \begin{bmatrix} \tau_1 \\ \tau_2 \end{bmatrix} = \begin{bmatrix} (m_1+m_2)a_1^2 + m_2a_2^2 + 2m_2a_1a_2c_2 & m_2a_2^2 + m_2a_1a_2c_2 \\ m_2a_2^2 + m_2a_1a_2c_2 & m_2a_2^2 \end{bmatrix} \begin{bmatrix} \ddot{\theta}_1 \\ \ddot{\theta}_2 \end{bmatrix} + \begin{bmatrix} -m_2a_1a_2s_2\dot{\theta}_2^2 - 2m_2a_1a_2s_2\dot{\theta}_1\dot{\theta}_2 \\ m_2a_1a_2s_2\dot{\theta}_1^2 \end{bmatrix} + \begin{bmatrix} (m_1+m_2)ga_1c_1 + m_2ga_2c_{12} \\ m_2ga_2c_{12} \end{bmatrix}
\end{equation}

这些在标准形式中
\begin{equation}
    M(q)\ddot{q} + N(q, \dot{q}) = \tau
\end{equation}

取连杆质量为1 kg，长度为1 m。

PD计算力矩控制律给出为
\begin{equation}
    \tau = M(q)(\ddot{q}_d + K_v\dot{e} + K_pe) + N(q, \dot{q})
\end{equation}
其中跟踪误差定义为
\begin{equation}
    e = q_d - q
\end{equation}

\textbf{b. 期望轨迹}

让期望轨迹$q_d(t)$的分量为
\begin{align}
    \theta_{1d} &= g_1\sin(2\pi t/T) \\
    \theta_{2d} &= g_2\cos(2\pi t/T)
\end{align}
周期$T = 2$ s，幅值$g_i = 0.1$ rad $\approx 6$ deg。为了良好的跟踪，选择闭环系统的时间常数为0.1 s。对于临界阻尼，这意味着$K_v = \mathrm{diag}\{k_v\}, K_P = \mathrm{diag}\{k_p\}$，其中
\begin{equation}
    k_p = \omega_n^2 = 100, \quad k_v = 2\omega_n = 20
\end{equation}

重要的是要意识到控制器参数（如PD增益）的选择取决于性能目标——在这种情况下，是期望轨迹的周期。

\textbf{c. 计算机仿真}

让我们使用附录B中的程序TRESP仿真计算力矩控制器。使用MATLAB和SIMNON等商业包进行仿真非常相似。

TRESP所需的子程序如图4.4.2所示。它们是
\end{examplebox}

\begin{examplebox}[示例4.4-2：PID计算力矩控制的仿真]
在示例4.4.1中我们仿真了双连杆平面臂的PD计算力矩控制器。在此示例中，我们向臂动力学添加恒定的未知扰动，并比较PD与PID计算力矩控制。

设臂动力学为
\begin{equation}
    M(q)\ddot{q} + N(q, \dot{q}) = \tau + \tau_d
\end{equation}
其中$\tau_d$是每个分量为1 N-m的恒定扰动。这可以模拟未知动力学，如摩擦等。1 N-m的值代表了相当大的偏置。

在图4.4.2的子程序arm(x, xp)中，向非线性项$N_1$和$N_2$的计算添加1，并使用PD计算力矩控制器，取$k_p=100$，$k_v=20$，得到图4.4.7(a)中的误差图。由于未建模的恒定扰动未被计算力矩定律考虑，跟踪误差中存在不可接受的残余偏置。最大误差为0.033 rad，略小于2度。

现在添加积分误差加权项，取$k_i=500$，得到图4.4.7(b)的结果，显示出零稳态误差，效果相当好。

相关的控制力矩如图4.4.8所示，表明使用积分项并不会显著增加力矩幅值。

为了仿真PID控制律，有必要向子程序arm(x, xp)中的$x$添加两个额外的状态。方便起见，称它们为$x(5)$和$x(6)$，因此添加到子程序中的行是
\begin{align}
    \dot{x}(5) &= e(1) \\
    \dot{x}(6) &= e(2)
\end{align}

因此，现在有必要在子程序arm(x, xp)中计算$e(1)$和$e(2)$以进行积分。
\end{examplebox}

\subsection{PID外环设计}

为了提高性能，可以使用PID外环控制
\begin{equation}
    u = -K_v\dot{e} - K_pe - K_i\int e\,\dd t
\end{equation}

然后整体机器人臂输入变为
\begin{equation}
    \tau_c = M(q)\left(\ddot{q}_d + K_v\dot{e} + K_pe + K_i\int e\,\dd t\right) + N(q, \dot{q})
\end{equation}

此控制器如图4.4.6所示。闭环误差动力学为
\begin{equation}
    \ddot{e} + K_v\dot{e} + K_pe + K_i\int e\,\dd t = w
\end{equation}
或以状态空间形式，通过定义
\begin{equation}
    \dot{\varepsilon} = e
\end{equation}
我们有
\begin{equation}
    \frac{\dd}{\dd t}\begin{bmatrix} \varepsilon \\ e \\ \dot{e} \end{bmatrix} = \begin{bmatrix} 0 & I & 0 \\ 0 & 0 & I \\ -K_i & -K_p & -K_v \end{bmatrix}\begin{bmatrix} \varepsilon \\ e \\ \dot{e} \end{bmatrix} + \begin{bmatrix} 0 \\ 0 \\ I \end{bmatrix}w
\end{equation}

积分项的作用是消除稳态误差。然而，积分项可能导致执行器饱和引起的积分器饱和问题。

\subsection{计算力矩类控制器}

有许多计算力矩控制器的变体。在本小节中，我们总结了一些重要的变体。

考虑近似计算力矩控制器
\begin{equation}
    \tau_c = \hat{M}(\ddot{q}_d - u) + \hat{N}
\end{equation}
其中$\hat{M}$和$\hat{N}$是$M$和$N$的近似值，可能在实际实现中使用以节省计算量。然后闭环系统变为
\begin{equation}
    (I - \Delta)\ddot{e} + (I - \Delta)(K_v\dot{e} + K_pe) + \delta = w
\end{equation}
其中
\begin{equation}
    \Delta = I - \hat{M}^{-1}M, \quad \delta = \hat{M}^{-1}(N - \hat{N})
\end{equation}
和干扰为
\begin{equation}
    w = \hat{M}^{-1}\tau_d
\end{equation}

如果使用精确计算力矩控制使得$\Delta = 0, \delta = 0$，这会简化为误差系统(4.4.10)。否则，误差系统由期望加速度和非线性项不匹配$\delta(t)$驱动。因此，跟踪误差永远不会精确为零。此外，辅助控制$u(t)$乘以$(I - \Delta)$，这可能导致一个非常困难的控制问题。

使用外环PD反馈，使得$u(t) = -K_v\dot{e} - K_pe$，得到误差系统
\begin{equation}
    \ddot{e} + (I - \Delta)K_v\dot{e} + (I - \Delta)K_pe + d = w
\end{equation}
其中$d(t) = \delta + \Delta\ddot{q}_d$。即使$K_v$和$K_p$被选择为左侧具有良好的稳定性，这种系统的行为也不是显而易见的。存在两种问题：首先是干扰项$d(t)$，其次是误差及其导数的函数$\Delta(K_v\dot{e} + K_pe)$。

\subsection{PD-Plus-Gravity控制器}

计算力矩族中一个有用的控制器是PD-plus-gravity控制器，当$\hat{M} = I, \hat{N} = G(q) - \ddot{q}_d$时产生，其中$G(q)$是操作臂动力学的重力项。然后，为$u(t)$选择PD反馈产生
\begin{equation}
    \tau_c = -K_v\dot{e} - K_pe + G(q)
\end{equation}

此控制律在[Arimoto and Miyazaki 1984]、[Schilling 1990]中进行了处理。它比精确计算力矩控制器实现起来简单得多。

当臂静止时，动力学(4.4.1)中唯一的非零项是重力$G(q)$、干扰$\tau_d$和控制力矩$\tau$。PD-gravity控制器$\tau_c$包含$G(q)$，因此我们应该期望对设定点跟踪有良好的性能，即当给定恒定的$q_d$使得$\ddot{q}_d = 0$时。下一个结果将这一点形式化。它依赖于Lyapunov证明（第1章），这种证明在本书中将始终有用，特别是利用表3.3.1中的斜对称性质。因此，理解此证明中的步骤非常重要。

\begin{theorem}
假设在臂动力学(4.4.1)中使用PD-gravity控制且$\tau_d = 0, \ddot{q}_d = 0$。则稳态跟踪误差$e = q_d - q$为零。
\end{theorem}

\begin{proof}
\textbf{1. 闭环系统}

忽略摩擦，机器人动力学由下式给出
\begin{equation}
    M(q)\ddot{q} + V_m(q, \dot{q})\dot{q} + G(q) = \tau
\end{equation}

当$\ddot{q}_d = 0$时，提出的控制律(4.4.49)产生闭环动力学
\begin{equation}
    M\ddot{q} + V_m\dot{q} + K_v\dot{q} + K_pe = 0
\end{equation}

\textbf{2. Lyapunov函数}

现在选择Lyapunov函数
\begin{equation}
    V = \frac{1}{2}\dot{q}^TM\dot{q} + \frac{1}{2}e^TK_pe
\end{equation}

并微分得到
\begin{equation}
    \dot{V} = \frac{1}{2}\dot{q}^T\dot{M}\dot{q} + \dot{q}^TM\ddot{q} + \dot{e}^TK_pe
\end{equation}

代入闭环动力学(2)产生
\begin{equation}
    \dot{V} = \frac{1}{2}\dot{q}^T(\dot{M} - 2V_m)\dot{q} - \dot{q}^TK_v\dot{q}
\end{equation}

现在，第一项的斜对称性给出
\begin{equation}
    \dot{V} = -\dot{q}^TK_v\dot{q} \leq 0
\end{equation}

状态为$x = [e^T \, \dot{q}^T]^T$，因此$V$是正定的但只是负半定的。因此，我们展示了Lyapunov意义上的稳定性，即误差和关节速度都是有界的。

\textbf{3. LaSalle扩展的渐近稳定性}

可以使用Barbalat引理和LaSalle扩展的变体[Slotine and Li 1991]（第2章）来演示系统的渐近稳定性。因此，有必要演示包含在集合$\dot{V} = 0$中的唯一不变集是原点。

由于$V$的下界为零且非正，因此$V$趋近于一个有限极限，可以写成
\begin{equation}
    V(\infty) - V(0) = \int_0^\infty \dot{V}\,\dd t < \infty
\end{equation}

我们现在调用Barbalat引理来表明$\dot{V}$趋于零。为此，有必要演示$\dot{V}$的一致连续性。我们看到
\begin{equation}
    \ddot{V} = -2\dot{q}^TK_v\ddot{q}
\end{equation}

所展示的稳定性表明$\dot{q}$和$q$是有界的，因此(2)和$M^{-1}(q)$的有界性（见表3.3.1）揭示$\ddot{q}$是有界的。因此，$\ddot{V}$是有界的，因此$\dot{V}$是一致连续的。这保证了根据Barbalat引理，$\dot{V}$趋于零。

现在很明显$\dot{q}$趋于零。还需要表明跟踪误差$e(t)$消失。注意当$\dot{q} = 0$时，(2)揭示
\begin{equation}
    K_pe = 0
\end{equation}

因此，非零的$e(t)$导致非零的$\ddot{q}$，因此$\dot{q} \neq 0$，这是一个矛盾。因此，包含在$\{x(t) | \dot{V} = 0\}$中的唯一不变集是$x(t) = 0$。这最终表明$e(t)$和$\dot{q}$都消失，并结束了证明。
\end{proof}

\begin{examplebox}[示例4.4-3：PD-Gravity控制器的仿真]
在示例4.4.1中，我们在双连杆平面操作臂上仿真了精确的计算力矩控制律。在这里，让我们仿真PD-gravity控制器。我们将采用相同的臂参数和期望轨迹。

假设每个连杆的PD增益相同，双连杆臂的PD-gravity控制律为
\begin{equation}
    \tau_c = -K_v\dot{e} - K_pe + G(q)
\end{equation}

这很容易通过注释掉几行代码并在图4.4.2的子程序CTL(x)中进行一些其他修改来仿真。

对于临界阻尼，PD增益选择为
\begin{equation}
    k_p = \omega_n^2, \quad k_v = 2\omega_n
\end{equation}

$\omega_n$几个值的PD-gravity控制器仿真结果如图4.4.9所示。随着$\omega_n$的增加，因此PD增益也增加，跟踪性能提高。然而，无论PD增益多大，跟踪误差永远不会精确为零，而是以零为中心的一个球内有界，其半径随着增益变大而减小。$\omega_n = 50$的性能，对应于$k_p = 2500, k_v = 100$，对于许多应用来说会相当合适。

非常重要的是要注意误差的直流值等于零。可以考虑将重力项作为计算力矩控制律的"直流部分"。如果将它们包括在计算力矩中，将没有误差偏移。

相关的控制输入力矩如图4.4.10所示。非常有趣的是，力矩对于更高的增益更小。这与流行的观点相反，后者假设控制力矩总是随着PD增益的增加而增加。这是由于更大的增益提供更好的性能，因此跟踪误差更小这一事实。

鉴于图4.4.9(a)中的误差具有不同的幅值，将内连杆的PD增益比外连杆更大可能更合理。
\end{examplebox}

\subsection{经典关节控制}

通过选择(4.4.43)中的
\begin{equation}
    \hat{M} = I, \quad \hat{N} = 0
\end{equation}
可以获得在实践中经常给出良好结果的简单控制器。

然后有结果
\begin{equation}
    \tau_c = -K_v\dot{e} - K_pe
\end{equation}

通过选择仅依赖于关节变量$i$的$u_i$，这描述了$n$个解耦的独立关节控制器，被称为独立关节控制。在实际机器人手臂上的实现很容易，因为关节$i$的控制输入仅依赖于局部测量的变量，而不依赖于其他关节的变量。此外，计算简单，不涉及求解复杂的非线性机器人逆动力学。

在机器人技术的早期，独立关节控制很受欢迎[Paul 1981]，因为它允许使用单输入/单输出(SISO)经典技术对闭环系统进行解耦分析。它也被称为经典关节控制。即使在今天，也有一些研究者强烈主张这种控制方案始终适合实际实现，并且现代非线性控制方案对于工业机器人应用来说过于复杂。

现在进行独立关节控制的传统分析。它提供了与经典控制概念的联系，对于机器人控制设计人员理解很重要。参见[Franklin et al. 1986]了解经典控制理论参考。

具有电动执行器的机器人臂的简化动力学模型可以写成（第3.6节）
\begin{equation}
    \left(J_{mi} + \frac{m_{ii}}{r_i^2}\right)\ddot{\theta}_i + \left(B_{mi} + \frac{k_{bmi}k_{mi}}{R_i}\right)\dot{\theta}_i + \frac{d_i}{r_i} = \frac{k_{mi}}{R_i}v_i
\end{equation}
其中$J_m$是执行器电机惯量，$B_m$是转子阻尼常数，$k_m$是力矩常数，$k_b$是反电动势常数，$R$是电枢电阻，$r_i$是关节$i$的齿轮比。电机角度表示为$\theta_i(t)$。$M(q)$的对角元素的常数部分表示为$m_{ii}$。这些元素的时间变化部分，以及$M(q)$的非对角元素、非线性项$N(q, \dot{q})$和任何干扰$\tau_d$都集中在一起成为干扰$d_i(t)$。因此$d_i(t)$包含所有其他关节对关节$i$的影响。控制输入是电机电枢电压$v_i(t)$。

请注意，在此方程中主要出现电机参数。事实上，如果齿轮比很小，甚至可以忽略$m_{ii}$。由于这个原因，如果齿轮比很小，机器人臂控制问题实际上就简化为控制执行器电机的问题。

让我们将操作臂关节$i$的这种简化线性时不变模型表示为
\begin{equation}
    J\ddot{\theta}_i + B\dot{\theta}_i + r d_i = k v_i
\end{equation}
其中常数$k_m/R$已被并入$J$和$B$的定义中。

根据表3.3.1中的性质，干扰$d(t)$是有界的，尽管通常不是由常数界定。界限可以是$q(t), \dot{q}(t)$甚至$\ddot{q}(t)$的函数。这通常在经典分析中被忽略。然而，$d(t)$的影响通过乘以齿轮比$r$而有所减轻。在经典关节控制设计中忽略了关节柔性（第6章）的影响。[为了与经典符号保持一致，本节中$u(t)$的定义与我们以前的用法相比有一个符号变化。]

现在，让我们考虑控制输入$u(t)$的一些选择。

\textbf{PD控制}。选择PD控制律
\begin{equation}
    u = k_v\dot{e} + k_pe
\end{equation}
其中$e(t) = \theta_{di}(t) - \theta_i(t)$是电机角度$i$的跟踪误差，得到图4.4.11所示的闭环系统。回想一下，电机角度与关节变量通过$q_i = r_i\theta_i$相关联。因此可以从$q_{di}(t)$计算$\theta_d(t)$。这个PD控制器很容易在实际的臂上实现，只需要很少的计算机能力。

使用经典控制中的Mason定理，发现设定点跟踪的闭环传递函数为（参见习题）
\begin{equation}
    \frac{\Theta_i}{\Theta_{di}} = \frac{k_ps + k_p}{Js^3 + (B + k_v)s^2 + k_vs + k_p}
\end{equation}
具有闭环特征多项式
\begin{equation}
    \Delta_c(s) = Js^3 + (B + k_v)s^2 + k_vs + k_p
\end{equation}
可以选择PD增益以获得合适的自然频率和阻尼比，正如我们所见。

在稳态时，对干扰$d(t)$的唯一非零贡献是被忽略的重力$G(q)$。表3.3.1显示，对于给定的机器人臂，重力向量由已知值$g_b$界定。因此，在稳态时，使用PD控制的关节$i$的稳态跟踪误差由
\begin{equation}
    |e_i| \leq \frac{r_i g_b}{k_p}
\end{equation}
界定。因此，对于指定最终值$q_d$的设定点跟踪，具有大$k_p$的PD控制可能是合适的。

事实上，下一个结果表明，即使对于跟随期望轨迹，而不仅仅是设定点控制，PD控制通常也非常合适。它在[Dawson 1990]中被证明。

\begin{theorem}
如果将PD控制律(4.4.54)应用于每个关节且$e(0) = 0, \dot{e}(0) = 0$，则位置和速度跟踪误差在球内有界，其半径随着$k_v \to \infty$大致以$1/\sqrt{k_v}$减小。
\end{theorem}

这个结果证实了那些坚持认为PD控制通常足够用于实际应用的人的观点。

当然，问题是$k_v$不能在没有达到执行器力矩限制的情况下无限增加。本书中将讨论的其他方案允许良好的轨迹跟随而不需要如此大的力矩。

在[Paul 1981]中讨论了修改PD控制律以获得更好性能的几种方法。这些包括重力补偿[这恰好产生控制器(4.4.49)]和加速度前馈，这相当于使用
\begin{equation}
    u = k_v\dot{e} + k_pe + \ddot{\theta}_d
\end{equation}
对于每个关节。还提到了关节耦合控制，这相当于在$u(t)$中重新添加$M(q)$和$N(q, \dot{q})$中描述关节之间相互作用的一些被忽略的项。因此，这种修正涉及在近似计算力矩控制(4.4.43)中使用更好的$\hat{M}$和$\hat{N}$估计。

\textbf{PID控制}。已经看到PD独立关节控制通常适合跟踪控制。然而，在稳态时存在由于重力引起的残余误差(4.4.57)。这可以使用PID独立关节控制律消除
\begin{equation}
    u = k_v\dot{e} + k_pe + k_i\int e\,\dd t
\end{equation}
对于每个关节。

现在设定点跟踪$\ddot{\theta}_d = 0$的传递函数为
\begin{equation}
    \frac{\Theta_i}{\Theta_{di}} = \frac{k_is + k_ps + k_i}{Js^3 + (B + k_v)s^2 + k_ps + k_i}
\end{equation}
具有闭环特征多项式
\begin{equation}
    \Delta_c(s) = Js^3 + (B + k_v)s^2 + k_ps + k_i
\end{equation}
现在终值定理显示，对于设定点控制，稳态误差为零。Routh检验表明，对于稳定性，需要
\begin{equation}
    k_i < \frac{(B + k_v)k_p}{J}
\end{equation}

在控制律中使用积分项时，重要的是要注意由于执行器饱和限制引起的积分器饱和的可能性，这在本节前面的"PID外环设计"中进行了讨论。

\begin{examplebox}[示例4.4-4：经典关节控制和力矩饱和限制]
在示例4.4.1中，我们为双连杆平面肘臂仿真了精确的计算力矩控制律。在这里，我们希望在同一臂上显示使用PD和PID经典关节控制的结果（具有相同的期望轨迹）。我们还对演示力矩饱和限制的影响感兴趣。为了突出效果而没有无关的细节，我们将只使用臂动力学，没有执行器动力学或齿轮比。

这些可以很容易地通过在图4.4.2中的子程序arm(x, xp)中进行一些轻微的修改来包括。

\textbf{a. PD独立关节控制}

对于双连杆臂，PD独立关节控制就是
\begin{equation}
    \tau_c = -K_v\dot{e} - K_pe
\end{equation}
其中跟踪误差$e(t) = q_d(t) - q(t)$。这个控制律非常简单直接，甚至不需要使用DSP就可以在实际的臂上实现。它比示例4.4.1中的控制简单得多。

在两连杆臂上使用PD控制器的结果如图4.4.13所示，对应于几个闭环自然频率$\omega_n$的值。回想一下，临界阻尼的PD增益为
\begin{equation}
    k_p = \omega_n^2, \quad k_v = 2\omega_n
\end{equation}

图(a)部分中的低增益结果($k_p = 100, k_v = 20$)非常糟糕（注意刻度）。然而，(c)部分中的高增益结果要好得多。更高的增益将导致进一步的跟踪误差减小。

重要的是要注意，在所有情况下，跟踪误差都有一个直流分量。这是由于重力项被忽略的事实，应该与示例4.4.3的结果形成对比。

因此在控制律中加入重力项显著提高了跟踪性能。

\textbf{b. PID独立关节控制}

双连杆臂的PID独立关节控制为
\begin{equation}
    \tau_c = -K_v\dot{e} - K_pe - K_i\int e\,\dd t
\end{equation}
这个简单的定律很容易通过向图4.4.2中的子程序arm(x, xp)添加两个状态$x(5)$和$x(6)$来实现，以考虑积分器（参见示例4.4.2）。

图4.4.13(c)中的仿真被重复，但现在添加$k_i = 1000$的积分增益。结果如图4.4.15所示。在正弦运动的几次迭代之后，跟踪误差的直流值变为零，因此性能大大提高。也就是说，添加积分项可以补偿控制律中忽略的重力项。然而，误差收敛到零直流值需要一段时间，结果仍然不如示例4.4.3中使用相同PD增益的PD-gravity控制器。另一方面，积分项可以拒绝除重力之外的其他项（例如，某些类型的摩擦）。

\textbf{c. 执行器饱和限制}

我们讨论了本例(b)部分前面"PID外环设计"中由于执行器饱和限制引起的积分器饱和。在这里，我们想演示其有害影响。

本例(b)部分中的控制力矩类似于图4.4.14(c)中的力矩；除了时间零附近的一些疯狂动作外，它表现得相当好，最大正偏移为78 N-m。

假设现在有力矩限制$\tau_{max} = 35$ N-m, $\tau_{min} = -35$ N-m。这很容易通过在图4.4.2中的子程序CTL(x)中添加几行代码来仿真
\begin{equation}
    \text{if } \tau(1) > 35 \text{ then } \tau(1) = 35
\end{equation}
等等。

这些限制的仿真结果如图4.4.16所示，非常糟糕。力矩如图4.4.17所示。

在4.5节中，我们展示了如何使用数字控制器中的抗饱和保护来改善执行器饱和的影响。PD控制的仿真结果与PID控制的结果相比不那么糟糕，因为如果没有积分器存在，饱和限制的影响较小。
\end{examplebox}

\section{数字机器人控制}

许多机器人控制方案是复杂的，涉及评估非线性项的大量计算。因此，它们在数字信号处理器(DSP)上作为数字控制律实现。某些控制方案可以通过使用简单的近似来以离散形式实现。

\subsection{基于近似连续控制器的设计}

机器人控制器设计通常按以下方式进行。首先，在连续时间中设计控制律以获得良好的性能。然后，使用某种近似技术将连续控制器数字化以在DSP上实现。

考虑连续时间控制器
\begin{equation}
    u_c(t) = -K_v\dot{e} - K_pe - K_i\int e\,\dd t
\end{equation}
这可以以传递函数形式表示为
\begin{equation}
    U_c(s) = -K(s)E(s)
\end{equation}
其中
\begin{equation}
    K(s) = K_vs + K_p + \frac{K_i}{s}
\end{equation}

为了获得数字控制器，我们可以使用Euler近似$s \approx (z - 1)/T$，其中$T$是采样周期。这产生
\begin{equation}
    K(z) = K_v\frac{z - 1}{T} + K_p + \frac{K_iT}{z - 1}
\end{equation}

然而，更常用的近似是双线性变换(BLT)
\begin{equation}
    s \approx \frac{2}{T}\frac{z - 1}{z + 1}
\end{equation}
这对应于使用梯形规则近似积分。在此映射下，稳定连续系统的极点被映射到稳定离散系统的极点。

\subsection{执行器饱和和积分器抗饱和补偿}

执行器饱和导致的积分器饱和是机器人控制器外环PID回路中可能发生的问题。在示例4.4.4中，我们看到了积分器饱和的有害影响。使用数字实现可以很容易地实现抗饱和保护[Åström and Wittenmark 1984]、[Lewis 1992]。

假设控制器以传递函数形式给出
\begin{equation}
    R(z^{-1})v_k = T(z^{-1})r_k - S(z^{-1})w_k
\end{equation}
其中$z^{-1}$在时间域中被解释为单位延迟$T$秒。

通过向两边添加$A_0(z^{-1})u_k$可以获得具有抗饱和补偿的调节器
\begin{equation}
    A_0(z^{-1})v_k = T(z^{-1})r_k - S(z^{-1})w_k
\end{equation}
\begin{equation}
    u_k = \mathrm{sat}(v_k)
\end{equation}
其中$\mathrm{sat}(\cdot)$是饱和函数。

当$A_0(z)$的所有$n$个极点都在原点时，会发生特殊情况。然后调节器显示死区行为；在$n$个时间步之后，其状态保持限制在依赖于$w_k$和$u_k$的值的容易计算的值（参见习题）。

\begin{examplebox}[示例4.5-1：数字机器人计算力矩控制器的仿真]
在示例4.4.1中，我们仿真了双连杆平面肘臂的连续时间PD计算力矩(CT)控制器。在此示例中，演示了该控制器离散化的一些问题。有两个目标。首先，我们展示了采样对CT控制器的影响。然后设计了一个仅使用编码器关节位置测量的实际数字控制器，重建计算控制律所需的速度。

\textbf{a. 数字CT控制器的采样周期影响}

在示例4.4.1中，给出了TRESP（附录B）所需的子程序，以仿真连续时间PD CT控制器。为了仿真PD CT控制器的数字化版本，有必要将控制器子程序[在该示例中称为CTL(x)]从连续动力学中移除并将其放入子程序DIG(IK, T, x)中，该子程序由TRESP调用。这与图4.5.1和4.3节中给出的数字控制器仿真技术保持一致。

图4.5.4中显示了TRESP所需的DIG(IK, TD, x)和F(t, x, xp)子程序。用于轨迹生成的子程序SYSINP(IT, x, t)和包含臂动力学的ARM(x, xp)与示例4.4.1中的相同。

现在使用不同的采样周期$T$运行程序TRESP得到图4.5.5所示的跟踪误差图。所有图的Runge-Kutta积分周期为5 ms。图表显示，对于$T = 5$ ms，误差非常小，实际上与示例4.4.1中的误差非常相似。然而，随着采样周期$T$的增加，跟踪性能恶化。

关节2的力矩输入如图4.5.6所示。对于小$T$，它实际上与示例4.4.1中的连续CT控制器相同。当$T = 0.1$ s时，会发生非常有趣的现象。根据图4.5.6(c)中的图，存在极限环，一种非线性振荡形式。众所周知，采样可以在闭环系统中引起这种非线性效应[Lewis 1992]。根据图4.5.5(d)，极限环在跟踪误差中反映为关于$e(t) = 0$的周期性振荡。极限环的出现与采样系统中可观测性的丧失密切相关。

\textbf{b. 仅使用位置测量的数字控制器}

在部分a中，我们假设关节位置和速度都可用作测量。在实际情况下，只有关节位置可用，通常来自光学编码器测量。因此，在这里，我们希望设计一个现实的数字CT控制器，重建速度。

图4.5.7中的子程序仅使用位置测量，使用导数滤波器(4.5.7)估计关节速度。产生的跟踪误差如图4.5.8所示。性能与部分a中使用速度测量的控制器相当，除了更大的初始误差瞬变。滤波器极点的值$v$取为0.1。

一些实现细节值得注意。首先，数字控制器的初始化方式非常重要。请注意在第一次迭代(IK=0)中将速度估计置零的代码行。删除这些行是一个很好的练习（参见习题）。

其次，可能会认为找到速度误差$\dot{e}_k = \dot{e}(kT)$的合理过程是找到关节速度$\dot{q}_k$的估计$v_k$，然后使用
\begin{equation}
    \dot{e}_k = \dot{q}_{dk} - v_k
\end{equation}
这有两个缺点。首先，它需要在内存中存储期望速度和期望轨迹。其次，它不起作用。

相反，有必要使用两个导数滤波器，一个用于估计关节速度，一个用于提供速度误差$\dot{e}_k$的估计。尝试使用$\dot{q}_{dk} - v_k$的仿真很好地说明了这一点。

请注意，速度估计也用于代替$\dot{q}$计算非线性项中的值。
\end{examplebox}

\begin{examplebox}[示例4.5-2：带抗饱和补偿的数字PI控制器]
采样周期为$T$秒的一般数字PI控制器由下式给出
\begin{equation}
    u_k = u_{k-1} + k_p e_k + \frac{k_p T}{T_I} \sum_{i=0}^{k} e_i
\end{equation}

我们使用改进的匹配极零点(MPZ)方法对连续PI控制器进行采样，在积分器中获得$T$秒的延迟以允许计算时间。比例增益为$k_p$，重置时间为$T_I$；两者在设计阶段都是固定的。跟踪误差为$e_k$。

乘以$(1-z^{-1})$写成
\begin{equation}
    (1-z^{-1})U(z) = k_p(1-z^{-1})E(z) + \frac{k_p T}{T_I} E(z)
\end{equation}

这是传递函数形式。用于实现的相应差分方程形式为
\begin{equation}
    u_k = u_{k-1} + k_p e_k + \frac{k_p T}{T_I} e_k
\end{equation}

由于自回归多项式$R=1-z^{-1}$在$z=1$处有根，使其处于边缘稳定，因此该控制器将经历饱和问题。因此，当$u_k$受限时，积分器将继续积分，"饱和"超出饱和水平。

\textbf{a. 抗饱和补偿}

为了纠正这个问题，选择一个观测器多项式
\begin{equation}
    A_0(z) = z - \alpha
\end{equation}

其中$|\alpha|<1$在期望位置有一个极点。设计参数$\alpha$可以通过仿真研究选择。然后具有抗饱和保护的控制器由下式给出
\begin{align}
    u_k &= u_{k-1} + k_p(e_k - e_{k-1}) + \frac{k_p T}{T_I} e_k + (1-\alpha)(u_{k-1} - u_{k-1}^{sat}) \\
    u_k^{sat} &= \mathrm{sat}(u_k)
\end{align}

如果$\alpha=1$，我们获得特殊情况，称为位置形式，没有抗饱和补偿。如果$\alpha=0$，我们获得死拍抗饱和补偿
\begin{equation}
    u_k = u_{k-1} + k_p e_k + \frac{k_p T}{T_I} e_k + (u_{k-1} - u_{k-1}^{sat})
\end{equation}

具有相应的差分方程实现。如果$u_k$未饱和，这相当于通过将(10)右侧的第二和第三项加到$u_{k-1}$来更新工厂控制。这些项因此不过是$u_k - u_{k-1}$。因此，$\alpha=0$的补偿器称为PI控制器的速度形式。

\textbf{b. 直流电机的数字控制}

考虑由下式给出的简化直流电机模型
\begin{equation}
    J\dot{\omega} + b\omega = \tau
\end{equation}

其中$\omega$是角速度。电机速度控制器具有形式
\begin{equation}
    \tau_k = k_p e_k + \frac{k_p T}{T_I} \sum_{i=0}^{k} e_i
\end{equation}

其中$e=r-\omega$是跟踪误差，$r$是期望的命令角速度。
\end{examplebox}

\section{最优外环设计}

在4.4节中，我们讨论了计算力矩控制，展示了如何使用涉及逆操作臂动力学的精确技术以及通过各种近似方法选择内控制环。我们还讨论了设计外线性反馈（跟踪）环的几种方案。本节讨论的结果总结在表4.4.1中。在本节中，我们打算介绍用于选择外反馈环的现代控制最优技术。现代最优设计在存在干扰和未建模动力学的情况下产生改进的鲁棒性。

\subsection{线性二次最优控制}

首先，有必要回顾现代线性二次(LQ)设计。假设我们给出了状态空间形式的线性时不变系统
\begin{equation}
    \dot{x} = Ax + Bu
\end{equation}
具有$n$维状态$x(t)$和$m$维输入$u(t)$。希望计算状态反馈增益$K$
\begin{equation}
    u = -Kx
\end{equation}
使得闭环系统
\begin{equation}
    \dot{x} = (A - BK)x
\end{equation}
是渐近稳定的。此外，我们不想使用太多控制能量来稳定系统，因为在许多现代系统（例如汽车、航天器）中，燃料或能量是有限的。

这是一个复杂的多变量设计问题，因为反馈增益矩阵$K$的维度为$m \times n$。经典控制方法可能涉及，例如，执行$mn$根轨迹设计以一次关闭一个反馈环。另一方面，现代控制技术通过求解一些标准矩阵设计方程可以找到保证稳定性的解决方案。这种方法同时关闭所有$nm$反馈环并保证良好的增益和相位裕度。

使用现代控制理论找到反馈矩阵如下。首先，定义形式为的二次性能指标(PI)
\begin{equation}
    J = \frac{1}{2}\int_0^\infty (x^TQx + u^TRu)\,\dd t
\end{equation}
其中$Q$是对称半正定$n \times n$矩阵（表示为$Q \geq 0$），$R$是对称正定$m \times m$矩阵（$R > 0$）。也就是说，$R$的所有特征值都大于零，$Q$的所有特征值都大于或等于零。$Q$被称为状态加权矩阵，$R$是控制加权矩阵。这些矩阵是设计参数，由工程师选择，例如，取决于闭环时间响应的期望形式。

最优LQ反馈增益$K$是使PI$J$最小化的增益。动机如下。PI中的二次项$x^TQx$和$u^TRu$是广义能量函数（例如，电容器中的能量为$\frac{1}{2}Cv^2$，运动的动能为$\frac{1}{2}mv^2$）。假设在闭环系统(4.6.3)中最小化$J$。这意味着$x^T(t)Qx(t) + u^T(t)Ru(t)$的无限积分是有限的，因此这个时间的函数随着$t$变大而趋于零。然而
\begin{equation}
    x^TQx = \|\sqrt{Q}x\|^2, \quad u^TRu = \|\sqrt{R}u\|^2
\end{equation}
其中矩阵的平方根定义为$\sqrt{Q}\sqrt{Q} = Q$。由于这些范数随着$t$消失且$|R| \neq 0$，函数$\|\sqrt{Q}x\|$和$u(t)$都趋于零。在假设$\sqrt{Q}, A$是可观测的[Kailath 1980]的情况下，如果$y(t)$趋于零，$x(t)$也趋于零。

因此，最优增益$K$保证在闭环系统(4.6.3)中所有信号都随时间趋于零。也就是说，$K$稳定$(A - BK)$。

在现代控制理论中，确定最优$K$很容易，是标准结果（参见，例如[Lewis 1986a]、[Lewis 1986b]）。只需通过求解矩阵设计方程就可以找到最优反馈增益
\begin{equation}
    K = R^{-1}B^TP
\end{equation}
\begin{equation}
    0 = A^TP + PA - PBR^{-1}B^TP + Q
\end{equation}
其中$P$是依赖于最优增益的辅助设计对称$n \times n$矩阵。第二个是被称为Riccati方程的非线性矩阵二次方程；使用MATLAB、[MATRIXx 1989]等软件设计包中的标准例程，很容易求解该方程以获得辅助矩阵$P$。

下一个结果在现代控制理论中非常重要，形式化了刚才给出的稳定性讨论。

\begin{theorem}
设$\sqrt{Q}, A$可观测且$(A, B)$可控。则：
\begin{enumerate}[label=(\alph*)]
\item 存在Riccati方程的唯一正定解$P$。
\item 闭环系统$(A - BK)$是渐近稳定的。
\item 闭环系统具有无限增益裕度和$60°$相位裕度。
\end{enumerate}
\end{theorem}

可控性在第1章中讨论过。可观测性大致意味着所有系统模式在PI中都有独立的影响，因此如果$J$有界，所有模式都独立地随$t$趋于零。验证这些属性很容易。如果可控性矩阵
\begin{equation}
    U = \begin{bmatrix} B & AB & \cdots & A^{n-1}B \end{bmatrix}
\end{equation}
具有满秩$n$，则系统是可控的。如果可观测性矩阵
\begin{equation}
    V = \begin{bmatrix} \sqrt{Q} \\ \sqrt{Q}A \\ \vdots \\ \sqrt{Q}A^{n-1} \end{bmatrix}
\end{equation}
具有满秩$n$，则系统$(A, C)$是可观测的。MATLAB，例如，为这些测试提供例程。因此，可以选择状态加权矩阵$Q$以满足可观测性要求。

该定理使这种现代设计方法非常强大。无论输入和状态的数量如何，只要假设成立，就可以找到稳定系统的反馈增益。

\subsection{线性二次计算力矩设计}

现在我们将这些结果应用于机器人操作臂动力学的控制
\begin{equation}
    \ddot{e} = u + w
\end{equation}
根据4.4节，计算力矩控制律
\begin{equation}
    \tau_c = M(q)(\ddot{q}_d - u) + N(q, \dot{q})
\end{equation}
产生误差系统
\begin{equation}
    \ddot{e} = u + M^{-1}\tau_d
\end{equation}
我们可以将其写成
\begin{equation}
    \dot{x} = Ax + Bu
\end{equation}
状态定义为
\begin{equation}
    x = \begin{bmatrix} e \\ \dot{e} \end{bmatrix}
\end{equation}

现在，选择外环PD反馈
\begin{equation}
    u = -Kx = -\begin{bmatrix} K_p & K_v \end{bmatrix}\begin{bmatrix} e \\ \dot{e} \end{bmatrix} = -K_pe - K_v\dot{e}
\end{equation}

为了找到稳定增益$K$，选择PI中的设计参数$Q$为
\begin{equation}
    Q = \begin{bmatrix} Q_p & 0 \\ 0 & Q_v \end{bmatrix}
\end{equation}
使得位置和速度误差被独立加权。然后，由于代表$n$个解耦牛顿定律（即双积分器）系统的简单$A$和$B$矩阵形式，Riccati方程的解很容易找到（参见习题）。使用该解在(4.6.5)中产生最优稳定增益的公式
\begin{equation}
    K_p = R^{-1/2}\sqrt{Q_p}, \quad K_v = \sqrt{2K_p}
\end{equation}

此LQ方法揭示了PD增益与某些设计参数$Q$和$R$之间的关系，这些参数决定了闭环系统中的总能量。特别要注意的是，闭环系统中$x(t)$和$u(t)$的相对幅值可以权衡。事实上，如果$R$相对于$Q_p$和$Q_v$较大，则PI(4.6.4)中的控制力被状态加权更重。然后最优控制将通过选择较小的控制增益来使$u(t)$更小；因此响应时间将增加。另一方面，选择较小的$R$将增加PD增益并使误差更快消失。

如果$Q_p, Q_v$和$R$是对角的，那么PD增益$K_p, K_v$也是对角的。具有非对角$Q_p, Q_v$和$R$的LQ方法提供了关节之间耦合的外反馈环的可能性，有时可以改善性能。LQ设计的另一个重要特征是定理中提到的保证鲁棒性。这在近似计算力矩设计中非常有用，其中
\begin{equation}
    \tau_c = \hat{M}(\ddot{q}_d - u) + \hat{N}
\end{equation}
其中$\hat{M}$和$\hat{N}$可以是$M(q)$和$N(q, \dot{q})$的简化版本。这种控制器的性能可以预期超过使用任意选择$K_p$和$K_v$设计的控制器。

重要的是要注意，这种LQ设计导致$e(t), \dot{e}(t)$和$u(t)$方面的最小闭环能量。然而，机器人手臂的实际控制输入为
\begin{equation}
    \tau = M(q)(\ddot{q}_d - u) + N(q, \dot{q})
\end{equation}

尽管使用这种方法没有最小化$\tau(t)$中的能量，但我们可以使用一些范数不等式来写入
\begin{equation}
    \|\tau\| \leq \|M(q)\|\|\ddot{q}_d\| + \|M(q)\|\|u\| + \|N(q, \dot{q})\|
\end{equation}
因此，保持小的$\|u(t)\|$可能预期会使$\|\tau\|$更小。

我们已经使用计算力矩（即反馈线性化）方法推导了LQ控制器。产生基于相同Riccati方程的设计的另一种方法是采用完整的非线性臂动力学并找到Hamilton-Jacobi-Bellman方程的近似（即时不变）解[Luo and Saridis 1985]；[Luo et al. 1986]。

\section{笛卡尔控制}

我们已经看到了如何使机器人操作臂跟踪期望的关节空间轨迹$q_d(t)$。然而，在任何实际应用中，机器人手臂的期望轨迹都是在工作空间或笛卡尔坐标中给出的。一系列重要的论文处理了分解运动操作臂控制[Whitney 1969]；[Luh et al. 1980]、[Wu and Paul 1982]。在那里，关节运动被分解到笛卡尔坐标中，其中指定了控制目标。

有几种笛卡尔机器人控制方法。例如，可以：
\begin{enumerate}
\item 在3.5节中使用笛卡尔动力学进行控制设计（参见习题）。
\item 使用逆运动学将期望的笛卡尔轨迹$y_d(t)$转换为关节空间轨迹$q_d(t)$。然后使用关节空间计算力矩控制方案表4.4.1。
\item 使用笛卡尔计算力矩控制。
\end{enumerate}

\subsection{笛卡尔计算力矩控制}

这种方法从关节空间动力学开始
\begin{equation}
    M(q)\ddot{q} + N(q, \dot{q}) = \tau
\end{equation}

笛卡尔误差定义为
\begin{equation}
    e_y = y_d - y
\end{equation}
其中$y_d(t)$是期望的笛卡尔轨迹，$y(t)$是末端执行器笛卡尔位置。

末端执行器笛卡尔位置可以用$y(t)$以多种形式指定，包括：$4 \times 4$臂$T$矩阵；使用欧拉角的6向量；或使用四元数的7向量。

笛卡尔误差可以指定为6向量
\begin{equation}
    e_y = \begin{bmatrix} e_p \\ e_o \end{bmatrix}
\end{equation}
其中$e_p(t)$是位置误差，$e_o(t)$是方向误差。

假设$y = h(q)$，其中$h(q)$是从$q(t)$到$y(t)$的变换。那么相关的Jacobian为$J = \partial h/\partial q$且$\dot{y} = J\dot{q}$。

相对于$e_y(t)$的计算力矩控制由下式给出
\begin{equation}
    \tau = M(q)J^{-1}(\ddot{y}_d - \dot{J}\dot{q} - u) + N(q, \dot{q})
\end{equation}
产生误差系统
\begin{equation}
    \ddot{e}_y = u + w
\end{equation}
具有干扰
\begin{equation}
    w = JM^{-1}\tau_d - \dot{J}\dot{q}
\end{equation}
我们称(4.7.7)为笛卡尔计算力矩控制律。

外环控制$u(t)$可以使用已经提到的用于关节空间计算力矩控制的任何技术进行选择（参见表4.4.1）。对于PD控制，完整的控制律为
\begin{equation}
    \tau_c = M(q)J^{-1}(\ddot{y}_d - \dot{J}\dot{q} + K_v\dot{e}_y + K_pe_y) + N(q, \dot{q})
\end{equation}

笛卡尔计算力矩控制的一个缺点是需要计算逆Jacobian。为了避免在每个采样周期求逆Jacobian，我们可以使用近似笛卡尔计算力矩控制器
\begin{equation}
    \tau_c = \hat{M}\hat{J}^{-1}(\ddot{y}_d - \dot{J}\dot{q} + K_v\dot{e}_y + K_pe_y) + \hat{N}
\end{equation}
通过设置$\hat{M}\hat{J}^{-1} = I, \hat{N} = G(q)$获得此控制律的特殊情况，产生笛卡尔PD-gravity控制器
\begin{equation}
    \tau_c = -K_v\dot{e}_y - K_pe_y + G(q)
\end{equation}

\subsection{笛卡尔误差计算}

可以从使用臂$T$矩阵表示的臂运动学测量的关节变量计算实际笛卡尔位置
\begin{equation}
    y = T = \begin{bmatrix} n & o & a & p \\ 0 & 0 & 0 & 1 \end{bmatrix}
\end{equation}
并且可以类似地将期望笛卡尔位置表示为$y_d = T_d$。然后可以如下计算适合计算力矩控制的笛卡尔位置误差和速度误差。定义
\begin{equation}
    \dot{y} = \begin{bmatrix} v \\ \omega \end{bmatrix}
\end{equation}
其中$v(t), v_d(t)$是实际和期望的线速度，$\omega(t), \omega_d(t)$是实际和期望的角速度。然后
\begin{equation}
    \dot{e}_y = \begin{bmatrix} \dot{e}_p \\ \dot{e}_o \end{bmatrix} = \begin{bmatrix} v_d - v \\ \omega_d - \omega \end{bmatrix}
\end{equation}
很容易计算，因为$\dot{q}(t)$可以使用$\dot{q} = J^{-1}\dot{y}$从测量的关节变量确定。

线性位置误差简单地由
\begin{equation}
    e_p = p_d - p
\end{equation}
给出。可以获得用于反馈目的的合适方向误差$e_o(t)$更困难，但可以定义如下。

分别表示$T$和$T_d$的旋转变换部分为$R(t), R_d(t)$。方向误差可以用将$R(t)$带入$R_d(t)$的Euler轴$k(t)$的$\phi(t)$弧度旋转表示。事实上，可以定义3向量
\begin{equation}
    e_o = k\sin\phi
\end{equation}
其中可以假设$e_o(t)$很小。使用此定义，可以证明$e_o(t)$从$T$和$T_d$使用
\begin{equation}
    e_o = \frac{1}{2}(n \times n_d + o \times o_d + a \times a_d)
\end{equation}
找到。整体笛卡尔误差现在由$e_y = [e_p^T \, e_o^T]^T$给出。

\begin{quote}
\textit{注：读者可以通过}\href{https://frankjie09.github.io/Robot-Manipulator-Control-CN/figures/Computed_Torque_Control.html}{\textbf{交互式可视化工具}}\textit{对比传统PID控制与计算力矩控制的性能差异，观察非线性补偿对轨迹跟踪精度的影响。}
\end{quote}

\section{总结}

在本章中，我们展示了如何生成定义机器人末端执行器运动的平滑轨迹，该轨迹通过一组指定的点。然后我们涵盖了重要的计算力矩控制器类，它包含许多类型的机器人控制算法。经典和现代控制算法都由该类描述，因此计算力矩提供了旧运动控制算法与本书其余部分更现代算法之间的桥梁。

作为计算力矩算法的特殊类型，我们提到了PD控制、PID控制、PD-plus-gravity、经典关节控制和数字控制。大多数机器人控制算法都是数字实现的，计算力矩为分析数字化和采样周期大小的影响提供了严格的框架。这是通过将数字控制视为近似计算力矩定律并研究误差系统来实现的。

我们展示了现代线性二次外环设计的一些方面，并以笛卡尔控制的讨论结束。

\newpage
\section*{参考文献}
\addcontentsline{toc}{section}{参考文献}

\begin{enumerate}[label={[\arabic*]}, leftmargin=*]
\item Arimoto, S., and F. Miyazaki, ``Stability and robustness of PID feedback control for robot manipulators of sensory capability,'' \textit{Proc. First Int. Symp.}, pp. 783--799, MIT, 1984.

\item Åström, K.J., and B. Wittenmark, \textit{Computer Controlled Systems}. Englewood Cliffs, NJ: Prentice Hall, 1984.

\item Bobrow, J.E., S. Dubowsky, and J.S. Gibson, ``On the optimal control of robotic manipulators with actuator constraints,'' \textit{Proc. Am. Control Conf.}, pp. 782--787, June 1983.

\item Chen, Y.-C., ``On the structure of the time-optimal controls for robotic manipulators,'' \textit{IEEE Trans. Autom. Control}, vol. 34, no. 1, pp. 115--116, Jan. 1989.

\item Dawson, D.M., ``Uncertainties in the control of robot manipulators,'' Ph.D. thesis, School of Electrical Engineering, Georgia Institute of Technology, Mar. 1990.

\item Franklin, G.F., J.D. Powell, and A. Emami-Naeini, \textit{Feedback Control of Dynamic Systems}. Reading, MA: Addison-Wesley, 1986.

\item Geering, H.P., L. Guzzella, S.A.R. Hepner, and C.H. Onder, ``Time-optimal motions of robots in assembly tasks,'' \textit{IEEE Trans. Autom. Control}, vol. AC-31, no. 6, pp. 512--518, June 1986.

\item Gilbert, E.G., and I.J. Ha, ``An approach to nonlinear feedback control with applications to robotics,'' \textit{IEEE Trans. Syst. Man Cybern.}, vol. SMC-14, no. 6, pp. 879--884, Nov./Dec. 1984.

\item Gourdeau, R., and H.M. Schwartz, ``Optimal control of a robot manipulator using a weighted time-energy cost function,'' \textit{Proc. IEEE Conf. Decision Control}, pp. 1628--1631, Dec. 1989.

\item Hunt, L.R., R. Su, and G. Meyer, ``Global transformations of nonlinear systems,'' \textit{IEEE Trans. Autom. Control}, vol. AC-28, no. 1, pp. 24--31, Jan. 1983.

\item Jayasuriya, S., and M.-S. Suh, ``Sub-optimal control strategies for manipulators with actuator constraints: the near minimum-time problem,'' \textit{Proc. Am. Control Conf.}, pp. 61--62, June 1985.

\item Johansson, R., ``Quadratic optimization of motion coordination and control,'' \textit{IEEE Trans. Autom. Control}, vol. 35, no. 11, pp. 1197--1208, Nov. 1990.

\item Kahn, M.E., and B. Roth, ``The near-minimum-time control of open-loop articulated kinematic chains,'' \textit{Trans. ASME J. Dyn. Syst. Meas. Control}, pp. 164--172, Sept. 1971.

\item Kailath, T., \textit{Linear Systems}. Englewood Cliffs, NJ: Prentice Hall, 1980.

\item Kim, B.K., and K.G. Shin, ``Suboptimal control of industrial manipulators with a weighted minimum time-fuel criterion,'' \textit{IEEE Trans. Autom. Control}, vol. AC-30, no. 1, pp. 1--10, Jan. 1985.

\item LaSalle, J., and S. Lefschetz, \textit{Stability by Liapunov's Direct Method}. New York: Academic Press, 1961.

\item Lee, C.S.G., and M.H. Chen, ``A suboptimal control design for mechanical manipulators,'' \textit{Proc. Am. Control Conf.}, pp. 1056--1061, June 1983.

\item Lewis, F.L., \textit{Optimal Control}. New York: Wiley, 1986 (a).

\item Lewis, F.L., \textit{Optimal Estimation}. New York: Wiley, 1986 (b).

\item Lewis, F.L., \textit{Applied Optimal Control and Estimation}. Englewood Cliffs, NJ: Prentice Hall, 1992.

\item Luh, J.Y.S., M.W. Walker, and R.P.C. Paul, ``Resolved-acceleration control of mechanical manipulators,'' \textit{IEEE Trans. Autom. Control}, vol. AC-25, no. 3, pp. 195--200, June 1980.

\item Luo, G.L., and G.N. Saridis, ``L-Q design of PID controllers for robot arms,'' \textit{IEEE J. Robot. Autom.}, vol. RA-1, no. 3, pp. 152--159, Sept. 1985.

\item Paul, R.P., \textit{Robot Manipulators}. Cambridge, MA: MIT Press, 1981.

\item Schilling, R.J., \textit{Fundamentals of Robotics}. Englewood Cliffs, NJ: Prentice Hall, 1990.

\item Shin, K.G., and N.D. McKay, ``Minimum-time control of robotic manipulators with geometric path constraints,'' \textit{IEEE Trans. Autom. Control}, vol. AC-30, no. 6, pp. 531--541, June 1985.

\item Slotine, J.-J.E., and W. Li, \textit{Applied Nonlinear Control}. Englewood Cliffs, NJ: Prentice Hall, 1991.

\item Whitney, D.E., ``Resolved motion rate control of manipulators and human prostheses,'' \textit{IEEE Trans. Man Machine Syst.}, vol. MMS-10, no. 2, pp. 47--53, June 1969.

\item Wu, C.-H., and R.P. Paul, ``Resolved motion force control of robot manipulator,'' \textit{IEEE Trans. Syst. Man Cybern.}, vol. SMC-12, no. 3, pp. 266--275, June 1982.
\end{enumerate}

\newpage
\section*{习题}
\addcontentsline{toc}{section}{习题}

\subsection*{4.2节}
\begin{enumerate}[label=4.2--\arabic*]
\item \textbf{最小时间控制}。当初始和最终速度不为零时，推导最小时间控制切换时间$t_s$[cf. (4.2.11)]。

\item \textbf{多项式路径插值}。希望将单个关节从$q(0)=0$, $\dot{q}(0)=0$通过点$q(1)=5$, $\dot{q}(1)=40$移动到最终位置/速度$q(2)=10$, $\dot{q}(2)=0$。确定此两个间隔路径所需的三次插值多项式。绘制生成的路径并验证它满足$q(t)$和$\dot{q}(t)$的指定要求。

\item \textbf{LFPB}。使用LFPB重复习题4.2--2。

\item \textbf{加速度匹配的多项式路径}。推导在通过点匹配位置、速度和加速度所需的插值多项式。
\end{enumerate}

\subsection*{4.3节}
\begin{enumerate}[label=4.3--\arabic*, resume]
\item \textbf{柔性耦合系统仿真}。使用计算机仿真重现示例3.6.1中带柔性耦合轴的电机的仿真结果。

\item \textbf{非线性系统仿真}。Van der Pol振荡器是一个具有一些有趣特性的非线性系统。仿真初始条件$x_1(0)=0.1, x_2(0)=0.1$下的动力学。使用参数值$\alpha=0.1$，然后$\alpha=0.8$。绘制$x_1(t)$和$x_2(t)$，以及相平面中的$x_2(t)$与$x_1(t)$。
\end{enumerate}

\subsection*{4.4节}
\begin{enumerate}[label=4.4--\arabic*, resume]
\item 证明(4.4.32)。

\item \textbf{PD计算力矩仿真}。使用各种PD增益值重复示例4.4.1。

\item \textbf{经典关节控制}。证明(4.4.55)、(4.4.57)、(4.4.60)和(4.4.62)。

\item \textbf{带负载不确定性的PD计算力矩}。在示例4.4.1中，假设$m_2$在$t=5$ s时从1 kg变为2 kg。CT控制器仍使用$m_2=1$的值。使用仿真绘制误差时间历史。

\item \textbf{带负载不确定性的PID计算力矩}。使用PID外环重复习题4.4--4。

\item \textbf{带摩擦不确定性的PD计算力矩}。将形式为$F(q, \dot{q}) = F_v\dot{q} + K_d\text{sgn}(\dot{q})$的摩擦添加到臂动力学中，但不添加到CT控制器中。仿真不同PD增益下的性能。

\item \textbf{带摩擦不确定性的PID计算力矩}。使用PID外环重复习题4.4--6。

\item \textbf{带执行器动力学的PD计算力矩}。
\begin{enumerate}[label=(\alph*)]
\item 为具有执行器动力学的双连杆平面肘臂设计CT控制律。取连杆质量和长度为1 kg, 1 m。取电机参数为$J_m=0.1$ kg-m$^2$，$b_m=0.2$ N-m/rad/s，$R=5\Omega$，设置齿轮比$r=0.1$。
\item 对各种PD增益值仿真控制器。
\end{enumerate}

\item \textbf{带忽略执行器动力学的PD计算力矩}。假设CT控制器仅使用臂动力学设计，忽略了执行器动力学。使用示例4.4.1中的PD增益。对各种齿轮比$r$值仿真控制律。

\item \textbf{仅使用执行器动力学的PD计算力矩}。仅使用执行器动力学设计CT控制器，不使用臂动力学。对各种齿轮比$r$值仿真控制律。

\item \textbf{带执行器动力学的经典关节控制}。像习题4.4--8那样包含执行器动力学重复示例4.4.4。

\item \textbf{带柔性耦合的PD计算力矩}。将示例3.6.1中的柔性轴与示例4.4.1中的双连杆臂结合起来。

\item \textbf{带近似CT控制的误差动力学}。考虑示例3.2.1中的双连杆极坐标臂，具有形式为$F(\dot{q})=F_v\dot{q}+K_d\text{sgn}(\dot{q})$的摩擦。对以下情况找到误差动力学(4.4.44)：
\begin{enumerate}[label=(\alph*)]
\item 摩擦未包含在CT控制律中。
\item 在CT控制律中负载质量$m_2$不完全已知。
\item 使用PD-gravity CT。
\item 使用PD经典关节控制，没有非线性项。
\end{enumerate}
\end{enumerate}

\subsection*{4.5节}
\begin{enumerate}[label=4.5--\arabic*, resume]
\item \textbf{数字控制仿真}。
\begin{enumerate}[label=(\alph*)]
\item 使用周期为1 s而不是2 s的期望轨迹重复示例4.5.1。
\item 通过删除图4.5.7中将初始速度估计置零的行重新进行仿真。
\item 尝试使用示例中方程(1)给出的从$\dot{q}_d$和$v_k$计算$\dot{e}_k$的替代技术仿真数字CT控制器。
\end{enumerate}

\item \textbf{数字控制仿真}。将示例4.4.3中的PD-gravity CT控制器转换为数字控制器。

\item \textbf{数字控制的误差动力学}。使用形式(4.5.3)的数字控制找到表4.4.1中的误差系统。

\item \textbf{抗饱和保护}。在示例4.4.4中的机器人控制器上实现抗饱和保护。
\end{enumerate}

\subsection*{4.6节}
\begin{enumerate}[label=4.6--\arabic*, resume]
\item \textbf{最优LQ外环PD增益}。验证(4.6.15)。

\item \textbf{使用LQ外环设计的鲁棒控制}。使用从LQ设计找到的PD增益重做示例4.4.3中的PD-gravity仿真。
\end{enumerate}

\subsection*{4.7节}
\begin{enumerate}[label=4.7--\arabic*, resume]
\item \textbf{直接笛卡尔计算力矩设计}。从3.5节中的笛卡尔动力学开始并设计计算力矩控制器。

\item \textbf{近似笛卡尔计算力矩}。推导与控制律(4.7.11)相关的误差系统动力学。

\item \textbf{笛卡尔PD-Plus-Gravity控制}。使用笛卡尔计算力矩控制重复示例4.4.3，其中轨迹在工作空间坐标中给出。
\end{enumerate}

\ifstandalone
  \end{document}
\fi